\section{Dyadic routing of a party network into a square-lattice PEPS}
\label{sec:peps_dyadic_routing}
This section treats the routing step of~\cite[Lemma~8.2]{OpenAI2026PolynomialPEPS}, September 24, 2026,
\texttt{07-assembly.tex}, lines 102--177: the links of a finite tensor network
whose parties sit at anchors of dyadic blocks are routed on the square grid with
constant edge congestion, folded into the genuine $L\times L$ square, and the
routed network is contracted into a PEPS on that square. The source derives
three properties of its party network (boundedly many parties per block,
bounded degree, link separation bounded by a fixed multiple of the endpoint
scales) from the distribution protocol
of~\cite[Proposition~7.1]{OpenAI2026PolynomialPEPS}, and a fourth (all anchors
in the padded square) from attaching each party to the anchor of its dyadic
block; its lemma has no hypotheses. Here these four properties are hypotheses
on an arbitrary party network, so the results below are a conditional form of
the lemma; the lemma for the network of the source follows once the protocol
and the attachment are shown to give these
properties~\cite{gap:polypeps_routing_network_bounds}.
% Status: that the protocol has these properties is not yet formalized.
The source also records that the two endpoint scales of every link
are equal or adjacent; the congestion bound below does not use this, so it is
not assumed.

\begin{definition}[Padded side and dyadic anchors]
    \label{def:peps_dyadic_anchor}
    \lean{TNLean.PEPS.Approximation.paddedSide, TNLean.PEPS.Approximation.dyadicAnchor}
    \leanok
    For $L\ge1$ let $N=2^{\lceil\log_2L\rceil}$ be the least power of two with
    $N\ge L$. A dyadic block of side $2^j$ with block index $r$ in one coordinate
    has anchor coordinate
    \[
        a_j(r)=\begin{cases}2^jr+2^{j-1},&j\ge1,\\ r,&j=0.\end{cases}
    \]
    A party attached to the block of scale $j$ and block indices $(r,s)$ sits at
    the anchor $(a_j(r),a_j(s))$.
\end{definition}

\begin{lemma}[Padded side]
    \label{lem:peps_dyadic_padded_side}
    \lean{TNLean.PEPS.Approximation.le_paddedSide, TNLean.PEPS.Approximation.paddedSide_lt_two_mul}
    \leanok
    \uses{def:peps_dyadic_anchor}
    For $L\ge1$ one has $L\le N<2L$.
\end{lemma}

\begin{proof}\leanok
    By minimality of the exponent, $2^{\lceil\log_2L\rceil-1}<L$ whenever the
    exponent is positive, so $N<2L$; the case $L=1$ gives $N=1$.
\end{proof}

\begin{lemma}[Anchor rows determine their scale]
    \label{lem:peps_dyadic_anchor_valuation}
    \lean{TNLean.PEPS.Approximation.padicValNat_two_dyadicAnchor, TNLean.PEPS.Approximation.dyadicAnchor_inj}
    \leanok
    \uses{def:peps_dyadic_anchor}
    For $j\ge1$ the two-adic valuation of $a_j(r)=2^{j-1}(2r+1)$ is exactly
    $j-1$. Consequently, if $j,j'\ge1$ and $a_j(r)=a_{j'}(r')$, then $j=j'$ and
    $r=r'$.
\end{lemma}

\begin{proof}\leanok
    The factor $2r+1$ is odd, so the valuation of $2^{j-1}(2r+1)$ is $j-1$.
    Equal anchors therefore have equal scales, and dividing by $2^j$ recovers the
    block index.
\end{proof}

\begin{definition}[Reflection fold]
    \label{def:peps_dyadic_fold}
    \lean{TNLean.PEPS.Approximation.foldCoord}
    \leanok
    On padded coordinates define $f(x)=x$ for $0\le x\le L-1$ and
    $f(x)=2(L-1)-x$ for $x\ge L$.
\end{definition}

\begin{lemma}[Properties of the fold]
    \label{lem:peps_dyadic_fold}
    \lean{TNLean.PEPS.Approximation.foldCoord_of_lt, TNLean.PEPS.Approximation.foldCoord_lt, TNLean.PEPS.Approximation.foldCoord_succ, TNLean.PEPS.Approximation.eq_or_eq_of_foldCoord_eq, TNLean.PEPS.Approximation.eq_or_eq_of_foldCoord_interval}
    \leanok
    \uses{def:peps_dyadic_fold}
    Let $L\ge1$. The fold fixes $\{0,\ldots,L-1\}$ and takes values there. For
    $x+2<2L$ the coordinates $f(x)$ and $f(x+1)$ are adjacent. For $x+1<2L$ every
    value $u=f(x)$ forces $x\in\{u,\,2(L-1)-u\}$, and for $x+2<2L$ every unit
    interval $\{u,u+1\}$ that is the image of $\{x,x+1\}$ forces
    $x\in\{u,\,2L-3-u\}$.
\end{lemma}

\begin{proof}\leanok
    Each statement is a case distinction on whether $x$ and $x+1$ lie below $L$,
    followed by linear arithmetic.
\end{proof}

\begin{definition}[Party network and routing]
    \label{def:peps_party_network_routing}
    \lean{TNLean.PEPS.Approximation.PartyNetwork, TNLean.PEPS.Approximation.PartyNetwork.coeff, TNLean.PEPS.Approximation.Routing}
    \leanok
    A party network with physical dimension $d$ consists of finite sets of
    parties and links, a dimension $D_\ell$ for each link $\ell$, two endpoint
    parties for each link, and for each party $p$ a tensor reading one label of
    every link incident to $p$ and one physical index. When the physical index of
    party $p$ is read at a site $\iota(p)$ of a finite simple graph $G$, the
    coefficient at a physical configuration $\sigma$ is the sum over all link
    labellings of the product of the party tensors evaluated at
    $\sigma(\iota(p))$. A routing of the network through $G$ assigns to every
    link a walk in $G$ from the site of its first party to the site of its second
    party.
\end{definition}

\begin{theorem}[Routed contraction]
    \label{thm:peps_routed_contraction}
    \lean{TNLean.PEPS.Approximation.Routing.toTensor_stateCoeff, TNLean.PEPS.Approximation.Routing.toTensor_bondDim_le, TNLean.PEPS.Approximation.Routing.exists_tensor_of_congestion_le}
    \leanok
    \uses{def:peps_party_network_routing}
    Let a party network be routed through a finite simple graph $G$. If every link
    has dimension at most $D$ and every edge of $G$ is traversed at most $\chi$
    times by the walks, counted with repetition, then there is a PEPS on $G$ with
    exactly the coefficients of the party network and every bond dimension at
    most ${\max(D,1)}^\chi$.
\end{theorem}

\begin{proof}\leanok
    Every step of a walk carries an identity tensor on the label of its link. At
    a vertex, the labels of all walks passing through it are kept independently,
    so crossings are permitted, and the party tensors located there are
    multiplied; zero-length links and coincident parties are contracted inside
    the vertex tensor. All traversals of one edge are grouped into a product
    alphabet, whose size is the product of the link dimensions, hence at most
    $D^m$ for $m$ traversals. Consistent configurations of the routed network are
    in bijection with link labellings, since the identity tensors force one label
    along each walk, so the contraction reproduces the coefficients of the party
    network.
\end{proof}

\begin{definition}[Horizontal-then-vertical routing on anchors]
    \label{def:peps_dyadic_hv_routing}
    \lean{TNLean.PEPS.Approximation.DyadicPlacement, TNLean.PEPS.Approximation.DyadicPlacement.anchor, TNLean.PEPS.Approximation.DyadicPlacement.linkRoute, TNLean.PEPS.Approximation.DyadicPlacement.load}
    \leanok
    \uses{def:peps_dyadic_anchor}
    Attach every party to a dyadic block, and so to its anchor. Each link,
    oriented arbitrarily, is routed from the anchor of its first party
    horizontally to the column of the anchor of its second party, then vertically
    to that anchor. The load of a unit edge of the padded square is the number of
    steps of all these walks that traverse it, counted with repetition.
\end{definition}

\begin{theorem}[Constant congestion on the padded square]
    \label{thm:peps_dyadic_padded_congestion}
    \lean{TNLean.PEPS.Approximation.DyadicPlacement.load_horizontal_le, TNLean.PEPS.Approximation.DyadicPlacement.load_vertical_le}
    \leanok
    \uses{def:peps_dyadic_hv_routing}
    Suppose at most $B$ parties are attached to any dyadic block, every party has
    at most $\Delta$ incident links, and the two anchors of every link differ in
    each coordinate by at most $c$ times the side of either endpoint block. Then
    every horizontal and every vertical unit edge of the padded square has load
    at most $2(2c+1)B\Delta$.
\end{theorem}

\begin{proof}\leanok
    \uses{lem:peps_dyadic_anchor_valuation}
    A horizontal step lies on the row of the anchor of the first party of its
    link, and each link traverses a given edge at most once. For a scale $j\ge1$
    the row $2^{j-1}(2s+1)$ determines $j$ and $s$ by its two-adic valuation.
    Anchors of one scale on one row are $2^j$ apart, so at most $2c+1$ blocks of
    that scale lie within horizontal distance $c2^j$ of the edge; scale zero gives
    a second bound of the same form. Each block carries at most $B$ parties, each
    the first endpoint of at most $\Delta$ links. Vertical steps are treated in
    the same way using the column of the second party.
\end{proof}

\begin{theorem}[Folded congestion]
    \label{thm:peps_dyadic_folded_congestion}
    \lean{TNLean.PEPS.Approximation.foldedRouting, TNLean.PEPS.Approximation.foldedRouting_card_traversal_le}
    \leanok
    \uses{thm:peps_dyadic_padded_congestion, lem:peps_dyadic_fold, lem:peps_dyadic_padded_side}
    Let $L\ge1$ and let every dyadic block lie in the padded square of side $N$.
    Under the hypotheses of the previous theorem, applying the fold in both
    coordinates turns the padded walks into walks of the open $L\times L$ square
    lattice, the party at an anchor is placed at its fold, and every edge of the
    genuine square is traversed at most $8(2c+1)B\Delta$ times, counted with
    repetition. This is the congestion part of~\cite[Lemma~8.2]{OpenAI2026PolynomialPEPS}
    for every party network satisfying the hypotheses stated above; the source
    states it for its own network without hypotheses.
\end{theorem}

\begin{proof}\leanok
    Since $N<2L$, every padded walk point has coordinates below $2L-1$, so
    adjacent padded points fold to adjacent genuine points. A horizontal edge of
    the genuine square has at most two padded preimages in each coordinate, hence
    at most four padded unit-edge preimages, and likewise for vertical edges.
    Summing the padded bound over these preimages gives
    $4\cdot2(2c+1)B\Delta$.
\end{proof}

\begin{theorem}[Dyadic routing yields a square-lattice PEPS]
    \label{thm:peps_dyadic_routing_peps}
    \lean{TNLean.PEPS.Approximation.exists_squarePEPS_of_dyadicRouting}
    \leanok
    \uses{thm:peps_dyadic_folded_congestion, thm:peps_routed_contraction}
    Let $L\ge1$ and consider a party network on dyadic anchors of the padded
    square with at most $B$ parties per block, at most $\Delta$ links per party,
    link separation at most $c$ times either endpoint scale, and link dimensions
    at most $D$. Then there is a PEPS on the open $L\times L$ square lattice whose
    coefficients equal those of the party network, with the physical index of
    each party read at the fold of its anchor, and whose bond dimensions are all
    at most ${\max(D,1)}^{\chi}$ with $\chi=8(2c+1)B\Delta$. The exponent $\chi$
    depends neither on $L$ nor on the number of dyadic scales.
\end{theorem}

\begin{proof}\leanok
    Apply the routed contraction to the folded routing, whose congestion is at
    most $\chi$ by the previous theorem. The fold fixes genuine coordinates, so a
    party whose anchor is a genuine site is read at that site.
\end{proof}
